\documentclass[11pt,letterpaper]{article}
\usepackage[margin=0.7in,headheight=15pt]{geometry}
\usepackage[T1]{fontenc}\usepackage{lmodern,amsmath,amssymb,fancyhdr,parskip,hyperref}
\hypersetup{colorlinks=true,urlcolor=blue}
\pdfinfoomitdate=1\pdftrailerid{}\pdfsuppressptexinfo=15
\pagestyle{fancy}\fancyhf{}\lhead{Week 80 / Midnight laboratory}\rhead{Facilitator draft}\cfoot{Bellingham Math Circle / Week 80 / facilitator / \thepage}
\setlength{\parskip}{7pt}\setlength{\parindent}{0pt}
\begin{document}
\section*{Mathematical overview}
Normalize midnight to time 1. The switch times $t_n=1-2^{-n}$ approach 1. With OFF initially and the post-switch convention, the successive states are ON, OFF, ON, OFF,\ldots. They have no limit: arbitrarily close to midnight both states occur. There is no last switch before midnight. Time convergence does not imply state convergence.

The instructions for $t<1$ alone do not determine the state at $t=1$. Appending ON or OFF as a separate endpoint assignment is mathematically consistent with those earlier instructions. Requiring continuity at 1 is impossible. Do not offer ``half on'' as an endpoint deduction or treat infinity as an even/odd integer.

By contrast, a counter starting at 0 and moving halfway toward 1 has positions $1-2^{-n}$ converging to 1. Experiments can show initial marks; the recurring halfway rule explains arbitrarily small remaining gaps. Grades 2--3 manipulate and compare; Grades 4--5 can separate an endpoint convention from a consequence. No actual supertask is performed.
\newpage
\section*{Preparation, readiness and first route}
Per pair: the printed time strip and a separate position strip, one counter, pencil and ruler. Cut the two ON/OFF panels and glue them back-to-back to make one reversible card. Work on the first four halves only, leaving a physically usable gap; describe or imagine later marks rather than folding them. Very late marks are mathematical descriptions, not demands for microscopic accurate folding. Put the two strips beside each other. Use the counter's center as its position; stop physical moves before counter size hides the gap. The endpoint remains visible and separate from switch marks.

Prerequisites: distinguish an endpoint from an interior mark; follow alternation; fold/locate halfway. Fraction notation is optional. One anchored adult reads the shared rule, witnesses each card flip and records a few predictions. The other child checks the move. No timer or flashing light is needed.

After the movement break, allow two minutes with strips/cards. In a three-minute whole-group launch, demonstrate one halfway mark and its card flip using a non-target instance; the worked visual gives input, action and output. Let children act before naming limits.

First route: for Grades 2--3 use Problems 1--2 only, with adult-scribed observations. Problems 3--5 are a Grades 3--5 continuation when endpoint discussion is ready. Allow 15 minutes marks and flips; 10 minutes halfway-counter motion on a fresh line; 10 minutes comparing what settles; five minutes considering extra endpoint rules. Preserve time for discussion. The rules are enforced by the partner checking marks and flipping the \emph{whole} two-state card. The optional endpoint conversation needs an adult; it is not a test of children saying the expected phrase. Timing and handling are untested.
\newpage
\section*{Answers and conversation moves}
Problem 1: a farther mark always exists: between any switch mark and midnight put the next halfway mark. Thus the ``last drawn mark'' is not a last mark of the rule. The original switch-time sequence has a visible endpoint in its description but no terminal event.

Problems 3--4: if a child predicts ON at midnight, invite an OFF group to identify a different additional assumption. Ask whether either is forced by the earlier moves. If a child says undefined, ask what is specified before midnight and what information would finish the instruction sheet. Preserve the distinction between a missing endpoint rule and a nonexistent ordinary limit.

Problems 2 and 5: for the moving counter, after each action the distance remaining is half the old distance. Given any positive real tolerance, sufficiently many halvings make the gap smaller. The counter never arrives in a finite move, but its positions have endpoint 1 as their limit. This is not a proof from six drawings: the rule continues for every finite number of actions.

For the lamp, the proposed settled-earlier-state rule in Problem 5 fails. Take the odd and even post-switch subsequences; one is always ON and the other always OFF. They cannot converge to one common state. Averages of states answer a different question and do not turn the binary lamp half on.

Hints belong here: try predicting the next \emph{two} states; keep the state card separate from the timeline; compare remaining distance with the previous distance. Do not provide those prompts after every small action on student pages.
\newpage
\section*{Encore, assumptions and sources}
Encore: design a process whose state becomes fixed after some finite mark, and another whose positions settle without arriving at a finite mark. Compare the number of actions, their times and their states separately. Return to the numbered-ball detective process: each label there does settle, even though the total count does not pass to the membership limit.

Assumptions: the two possible states are distinct; the rule covers all prescribed times strictly before midnight; a state immediately after a switch is used when describing the alternating sequence. Endpoint ON/OFF assignments do not claim physical realizability or continuity. No experiment with an actual infinitely fast mechanism is attempted.

Sources: James F. Thomson, ``Tasks and Super-Tasks'' (1954), especially the lamp discussion pp. 5--6; Paul Benacerraf, ``Tasks, Super-Tasks, and the Modern Eleatics'' (1962), pp. 768--770, distinguishing earlier behavior and endpoint specification. Original classroom presentation based on the organizer's example. Links and bibliographic details are in SOURCES.md.

Pedagogy: Rozhkovskaya, Lesson 7 ``At the lesson,'' item 1 (part0017) notes missed legal moves during game play. Partner witnessing is our response, informed also by the organizer's Week 2 experience. This lamp activity uses folding and an actual reversible card; it is not the old paper network-toggle rule.

Unpiloted. US Letter, single-sided, Actual Size. Physical fit, folding tolerances, timing and classroom rehearsal untested. After use record actual marks attempted, children's endpoint proposals, whether they distinguished time/state, and where the adult supplied the investigation rather than read/recorded.
\end{document}
